\chapter{Electronics}
\label{ch:electronics}

\section{Overview}

All electronic components are soldered onto a \qtyproduct{60 x 80}{mm} prototype board
(\cref{fig:pcb}). The main components are:
\begin{itemize}
  \item an \textbf{ESP32 development board} (HW-394, ESP32-32D module, USB-C) as the
        microcontroller,
  \item an \textbf{IMU breakout board} labelled \enquote{MPU-9250/6500/9255} (MPU-6500 family,
        three-axis accelerometer and gyroscope), connected via I\textsuperscript{2}C and read with
        the MPU6050\_light library~\cite{mpu6050light}, whose register accesses are compatible
        with this sensor,
  \item three \textbf{Nidec~24H brushless motors} with integrated driver and encoder; the
        \textbf{8-pin version} is required (the author first bought the 12-pin version, which does
        not fit this design),
  \item a \textbf{3S LiPo battery} (\SI{11.1}{V} nominal, \SI{12.6}{V} fully charged,
        \SI{1000}{mAh}) with an XT30 connector,
  \item a \textbf{mini DC-DC buck converter} (input \SIrange{5}{30}{V}, output \SI{5}{V}),
  \item an \textbf{active \SI{5}{V} buzzer} switched by a 2N2222 transistor,
  \item a \textbf{voltage divider} for measuring the battery voltage, and buffer capacitors.
\end{itemize}
The board carries three 8-pin headers for the motor cables and a 4-pin header for the IMU.
The complete schematic is shown in \cref{fig:schematic}; all connections are also listed in the
file \nolinkurl{hardware/electronics/wiring.md} of the repository.

The schematic was the author's first project in KiCad and is therefore not as clearly arranged as
it could be. Two details differ from the built cube: the transistor symbol is labelled 2N2219,
while a 2N2222 is used, and the motors are numbered 1--3 instead of X, Y, Z
(Nidec24H1~=~motor~X, Nidec24H2~=~motor~Z, Nidec24H3~=~motor~Y).

\begin{landscape}
\begin{figure}[p]
  \centering
  \includegraphics[width=\linewidth,height=0.88\textheight,keepaspectratio]{cubilox_mk1_schematic.pdf}
  \caption{Complete schematic of the Cubilox Mk1 (KiCad). The transistor Q1 is a 2N2222 in the
    built cube.}
  \label{fig:schematic}
\end{figure}
\end{landscape}

\section{Power supply}

The battery feeds a common \SI{12}{V} rail, which supplies the three motors directly (red wire,
VIN) and the buck converter. The buck converter produces \SI{5}{V} for the VIN pin of the ESP32
board and for the logic supply of the motors (white wire). The ESP32 board generates the
\SI{3.3}{V} for the IMU with its on-board regulator. All grounds are connected to one common GND
rail.

\begin{table}[htbp]
  \centering
  \caption{Power distribution.}
  \label{tab:power}
  \begin{tabularx}{\textwidth}{lX}
    \toprule
    Rail & Consumers \\
    \midrule
    Battery (\SIrange{11.1}{12.6}{V}) & motor power (VIN) of all three motors, buck converter input,
                                        voltage divider \\
    \SI{5}{V} (buck converter)        & ESP32 VIN, motor logic supply, buzzer \\
    \SI{3.3}{V} (ESP32 board)         & IMU \\
    \bottomrule
  \end{tabularx}
\end{table}

\subsection{Buffer capacitors}

The motors draw short, high current peaks whenever they accelerate or reverse, which happens
continuously while balancing. Because of the resistance and inductance of the battery and the
wires, these peaks cause short voltage dips on the supply rails. Such dips can disturb or reset
the ESP32 and the motor drivers. Three types of capacitors reduce this effect:
\begin{itemize}
  \item \textbf{C4, \SI{1000}{\micro F} electrolytic}, directly at the battery input: a large
        reservoir that supplies the current peaks of all motors together and smooths the main
        rail.
  \item \textbf{C1--C3, \SI{100}{\micro F} electrolytic}, one at the supply of each motor: local
        reservoirs close to the consumers, so that the peak current of one motor does not have to
        flow through the whole wiring and disturb the others.
  \item \textbf{C5, \SI{100}{nF} ceramic}, at the measuring point of the voltage divider: a small
        capacitor that filters noise on the ADC input (see \cref{sec:divider}).
\end{itemize}

\section{Battery voltage measurement}
\label{sec:divider}

The ESP32's analogue-to-digital converter (ADC) can only measure voltages up to about
\SI{3.3}{V}, while the battery reaches \SI{12.6}{V}. A voltage divider scales the battery voltage
down:
\begin{equation}
  V_\text{pin} = V_\text{bat}\,\frac{R_3}{R_2 + R_3}
             = V_\text{bat}\,\frac{\SI{10}{k\ohm}}{\SI{33}{k\ohm} + \SI{10}{k\ohm}}
             \approx 0.233\,V_\text{bat}.
\end{equation}
For a full battery, $V_\text{pin} = \SI{12.6}{V} \times 0.233 \approx \SI{2.93}{V}$, safely below
the limit. The firmware reverses this calculation by multiplying the measured pin voltage by
$(R_2+R_3)/R_3 = 4.3$ (\cref{lst:battery}). The upper resistor of \SI{33}{k\ohm} is built from
three resistors in series (\SI{20}{k\ohm} + \SI{10}{k\ohm} + \SI{3.3}{k\ohm}), because no single
\SI{33}{k\ohm} resistor was at hand.

The high resistance values keep the current through the divider small,
$I = \SI{12.6}{V} / \SI{43}{k\ohm} \approx \SI{0.29}{mA}$, so the divider hardly drains the battery.
Together with the source resistance of the divider,
$R_2 \parallel R_3 \approx \SI{7.7}{k\ohm}$, the \SI{100}{nF} capacitor C5 forms a low-pass filter
with a cut-off frequency of
\begin{equation}
  f_c = \frac{1}{2\pi\,(R_2\parallel R_3)\,C_5} \approx \SI{207}{Hz},
\end{equation}
which suppresses switching noise from the motors and provides the charge that the ADC needs for a
sample.

\begin{lstlisting}[caption={Battery measurement in \cmd{hardware.ino}: 16 samples are averaged and
  converted to the battery voltage.}, label={lst:battery}]
float readBatteryPinV() {
  uint32_t sum = 0;
  for (int i = 0; i < 16; i++) sum += analogReadMilliVolts(VBAT_PIN);
  return sum / 16.0 / 1000.0;
}

void updateBattery() {
  float v = readBatteryPinV() * VBAT_FACTOR;      // VBAT_FACTOR = (33 + 10) / 10
  batteryV = (batteryV <= 0) ? v : 0.8 * batteryV + 0.2 * v;
}
\end{lstlisting}

The battery voltage is only measured while the cube is not balancing, because the voltage drops
under load. Below \SI{11.1}{V} (\SI{3.7}{V} per cell) the cube warns with three beeps, and below
\SI{10.5}{V} (\SI{3.5}{V} per cell) it refuses to start, to protect the battery. A fully charged
battery noticeably increases the available motor power.

\section{Buzzer driver}

The active buzzer only needs to be switched on and off. Since an ESP32 pin can only supply a
small current, the buzzer is switched by an NPN transistor (2N2222) on its low side: GPIO~23 drives
the base through a \SI{1}{k\ohm} resistor, the emitter is connected to GND and the collector to
the negative terminal of the buzzer, whose positive terminal is connected to \SI{5}{V}. The base
current is approximately $(\SI{3.3}{V} - \SI{0.7}{V}) / \SI{1}{k\ohm} \approx \SI{2.6}{mA}$,
enough to switch the transistor fully on.

\section{Pin assignment}

\Cref{tab:pins} lists the pin assignment as defined in \cmd{config.h}. Strapping pins
(GPIO 0, 2, 5, 12, 15), the flash pins (GPIO 6--11) and the USB serial pins (GPIO 1, 3) were
deliberately avoided. GPIO 34, 35, 36 and 39 are input-only pins without internal pull
resistors; the encoders work on them without external pull-ups.

\begin{table}[htbp]
  \centering
  \caption{ESP32 pin assignment.}
  \label{tab:pins}
  \begin{tabular}{lcccc}
    \toprule
    Signal & Motor X & Motor Y & Motor Z & Other \\
    \midrule
    PWM        & 32 & 14 & 25 & \\
    Direction  & 4  & 13 & 27 & \\
    Encoder A  & 35 & 36 & 34 & \\
    Encoder B  & 33 & 19 & 18 & \\
    Brake (shared)        & \multicolumn{3}{c}{26} & \\
    I\textsuperscript{2}C SDA / SCL (IMU) & & & & 21 / 22 \\
    Buzzer                & & & & 23 \\
    Battery voltage (ADC) & & & & 39 \\
    \bottomrule
  \end{tabular}
\end{table}

\section{Motor connection}

The Nidec~24H contains its own driver: it only needs the supply voltages, a PWM signal for the
speed, a direction signal and a brake signal, and it returns two encoder signals.
\Cref{tab:motor_pins} shows the 8-pin connector. The colours of the wires can differ between
suppliers; the pin order is what matters. Two details are important for the firmware:
\begin{itemize}
  \item the PWM input is inverted: a duty cycle of 255 (of 255) means \enquote{stop}, so the
        firmware writes $255 - |u|$,
  \item the brake input is active low: HIGH lets the motors run freely, LOW brakes them.
\end{itemize}

\begin{table}[htbp]
  \centering
  \caption{Nidec~24H motor connector (identical for all three motors).}
  \label{tab:motor_pins}
  \begin{tabular}{clll}
    \toprule
    Pin & Colour & Function & Connected to \\
    \midrule
    1 & red    & VIN (battery voltage) & \SI{12}{V} rail \\
    2 & black  & GND                   & GND rail \\
    3 & yellow & brake                 & GPIO 26 (shared) \\
    4 & green  & PWM                   & motor-specific GPIO \\
    5 & blue   & direction             & motor-specific GPIO \\
    6 & white  & \SI{5}{V} logic supply & buck converter output \\
    7 & orange & encoder A             & motor-specific GPIO \\
    8 & brown  & encoder B             & motor-specific GPIO \\
    \bottomrule
  \end{tabular}
\end{table}

\section{Construction of the board}

\Cref{fig:pcb} shows both sides of the prototype board. The wiring on the bottom side is largely
colour-coded: red for supply voltages, black for GND, blue for the two I\textsuperscript{2}C lines
of the IMU, yellow for the brake and the buzzer, and green for all encoder and PWM signals.

\begin{figure}[htbp]
  \centering
  \begin{subfigure}[b]{0.49\textwidth}
    \centering
    \includegraphics[width=\textwidth]{pcb_top_view_annotated.jpg}
    \caption{Top side with the resistor values marked.}
    \label{fig:pcb_top}
  \end{subfigure}\hfill
  \begin{subfigure}[b]{0.49\textwidth}
    \centering
    \includegraphics[width=\textwidth]{pcb_bottom_view_wiring.jpg}
    \caption{Bottom side with the colour-coded wiring.}
    \label{fig:pcb_bottom}
  \end{subfigure}
  \caption{The prototype board: ESP32 board (right), buck converter (left), buzzer, voltage
    divider, buffer capacitors and the pin headers for the motors and the IMU.}
  \label{fig:pcb}
\end{figure}

The motor cables were extended with ordinary jumper wires with \SI{2.54}{mm} female connectors,
soldered to the motor wires and insulated with heat shrink tubing (\cref{fig:wiring_extension}).
This takes some time but is cheap and works well. Alternatives are ready-made extension cables
or crimping 8-pin cables with a matching connector kit.

\begin{figure}[htbp]
  \centering
  \includegraphics[width=0.45\textwidth]{soldered_wiring_extension.jpg}
  \caption{Motor cable extended with jumper wires, insulated with heat shrink tubing.}
  \label{fig:wiring_extension}
\end{figure}

\section{Safety note on the battery connection}

The XT30 battery connector is soldered to the board via pin headers. These contacts are not
insulated. They were useful during testing for measuring the battery voltage with a multimeter,
but they are dangerous: touching them or accidentally bridging them causes a short circuit of the
LiPo battery. \textbf{Do not copy this detail}; insulate all battery contacts.
